\documentclass[12pt]{article}
\usepackage{amsmath, amssymb, amsthm, geometry}
\geometry{a4paper, margin=1in}

\newtheorem{theorem}{Theorem}
\newtheorem{lemma}[theorem]{Lemma}
\newtheorem{proposition}[theorem]{Proposition}
\theoremstyle{definition}
\newtheorem{definition}[theorem]{Definition}
\newtheorem{remark}[theorem]{Remark}

\newcommand{\R}{\mathbb{R}}
\newcommand{\T}{\mathcal{T}}
\newcommand{\Q}{\mathcal{Q}}

\begin{document}

\title{Algebraic Characterization of Scaled Generic Quadrifocal Tensors}
\author{}
\date{}
\maketitle

\begin{abstract}
We provide a rigorous affirmative answer to the question of whether there exists a polynomial map characterizing the set of rank-1 scaled quadrifocal tensors arising from generic camera configurations. We construct such a map  explicitly using a combination of permutation consistency conditions and weighted Plücker relations. We prove that this map satisfies the required properties: it is independent of the specific camera matrices, has degrees independent of the number of views , and vanishes if and only if the scaling tensor is of rank 1.
\end{abstract}

\section{Introduction}

Let . Let  be Zariski-generic matrices representing  perspective cameras. For distinct indices , the quadrifocal tensor  is defined by:

We consider the set of tensors  defined by , where  is a scaling tensor supported on distinct indices. We seek a polynomial map  that vanishes precisely when  has rank 1 (i.e., ) on its support.

\section{Main Result}

\begin{theorem}
There exists a polynomial map  satisfying the following three properties:
\begin{enumerate}
\item  does not depend on the matrices .
\item The degrees of the coordinate functions of  do not depend on .
\item Let  satisfy  for precisely distinct . Then  if and only if there exist  such that  for all distinct .
\end{enumerate}
\end{theorem}

\section{Construction of the Map }

The variety of valid quadrifocal tensors is the intersection of the Grassmannian  (under the Plücker embedding) with the coordinate subspace defined by the camera geometry. The map  is constructed as the union of two sets of polynomial conditions: \emph{Permutation Consistency Constraints} and \emph{Weighted Plücker Constraints}.

\subsection{Permutation Consistency Constraints}
The generic tensors  satisfy strict antisymmetry under permutation of the camera indices. For example, swapping the first two cameras swaps the first two blocks of rows in the determinant, yielding:

For the scaled tensor , this symmetry is broken unless the scaling factors satisfy specific ratios. If  is rank-1, i.e., , we have:

This implies that for any fixed tuple of cameras, the tensors  and  must be proportional (with a coefficient independent of the spatial indices ). To enforce the global consistency of these ratios with a rank-1 structure, we impose cycle conditions.
We define polynomials  of degree 2 and higher. For every pair of tuples related by a swap of two indices, say , we require:

Furthermore, for every cycle of indices of length 3 (e.g., ), the product of the induced ratios must be consistent (equal to 1 or -1 depending on signature). These conditions depend only on subsets of indices of size .

\subsection{Weighted Plücker Constraints}
The unscaled tensors  satisfy the algebraic relations of the multiview variety. These are generated by the generalized Plücker relations on the  matrix of all camera rows. A standard Plücker relation is a homogeneous polynomial , where each  is a monomial in the entries of  (typically quadratic).
For the scaled tensor , simple substitution yields terms scaled by different factors. For instance, a quadratic relation involving views  might mix a term  (scaled by ) with  (scaled by ). For generic rank-1 , these scaling factors are not equal.

However, we can exploit the redundancy of the rank-1 structure. Since , we can construct \emph{balancing monomials}. For each term  in a Plücker relation, we multiply by a monomial  consisting of entries from other tensors in the collection such that the total scaling factor becomes identical for all terms in the sum.
Specifically, if term  has scaling factor , we choose  such that  is invariant. This is possible because the permutation group acts transitively on the rank-1 components. For example, the mismatch between  and  can be corrected by multiplying by ratios derived from  and  (as verified in the permutation step).
Thus, for every generator  of the multiview ideal (which involves at most 6 cameras), we define a polynomial:

\section{Proof of Properties}

\subsection{Independence and Bounded Degree}
The map  is defined by the structural equations of the permutation group and the Plücker relations of the Grassmannian . These are universal algebraic truths independent of the specific numerical values in . The relations are generated by local constraints involving subsets of cameras of size at most 6. Thus, the coordinate functions of  have a fixed maximum degree independent of .

\subsection{The Rank-1 Condition}
 Suppose  is rank-1. By construction,  satisfies the Permutation Consistency Constraints. For the Weighted Plücker Constraints, the balancing monomials ensure that the total scaling factor  is common to all terms. We have:

Thus, .

 Suppose . The Permutation Consistency Constraints imply that the ratios of  under permutation are consistent with a rank-1 tensor  (up to a scalar factor per tensor). Let  be the tensor such that . The Weighted Plücker Constraints then factor as . Since  is non-zero on the support, , implying . Thus,  satisfies the defining relations of the multiview variety. Since the Zariski-generic tensors are dense in the variety, and the variety is irreducible (for ),  corresponds to a valid camera configuration. Therefore,  is a rank-1 scaled quadrifocal tensor.

\end{document}
